% V2第19章思考题19.1--19.11，原扫描PDF185--186；图19.21--19.24需核验。
\section*{第19章\quad 磁场中的磁介质：思考题}
\paragraph{思考题 19.1}
判断这些磁场强度说法是否正确：(1) H只由自由传导电流决定；(2) 抗磁质和顺磁质中B总与H同向；(3) 对任何闭合曲面B的净通量为零；(4) 对任何闭合曲面H的净通量也为零。

\paragraph{思考题 19.2}
原图19.21把样品悬吊在通电螺线管的上部开口附近，若样品受向上的磁力，则它更可能为顺磁质还是抗磁质？说明场梯度条件。

\paragraph{思考题 19.3}
一闭合曲面只包围永磁体北极附近的材料而不包住南极，穿过该面的B总磁通与H总磁通分别怎样？

\paragraph{思考题 19.4}
永久磁铁落地可能部分退磁；把条磁铁朝南北方向放置、敲击又可能磁化，解释原因。

\paragraph{思考题 19.5}
铁磁体为何能吸引未被预先磁化的铁块？

\paragraph{思考题 19.6}
马蹄形磁铁不用时为何可在两极之间放置衔铁？不用的条磁铁又为何常N、S两极交错并列存放？

\paragraph{思考题 19.7}
铁磁质的磁导率强烈随温度变化，而许多抗磁质的磁化率与温度关系较弱，为什么？

\paragraph{思考题 19.8}
为什么磁路中的磁通量满足近似的串联不变、并联分支磁通相加规律？从磁感应强度散度解释。

\paragraph{思考题 19.9}
解释绝热去磁冷却为何要求顺磁盐先在低温磁化再与热源隔离移去外磁场；其适用的温度范围。

\paragraph{思考题 19.10}
中国古书中以磁针指向辨别方向的叙述蕴含了什么磁性认识？指南针的方向与地磁南北磁极、磁偏角之间怎样区分？

\paragraph{思考题 19.11}
图19.24(a)(b)电磁铁气隙窄或宽：气隙B与铁芯B是否完全相等？若安匝数NI相同，两种气隙B是否必然相同？
